\documentclass[11pt]{article}

\usepackage[a4paper,margin=28mm]{geometry}
\usepackage[T1]{fontenc}
\usepackage[utf8]{inputenc}
\usepackage{lmodern}
\usepackage{microtype}
\usepackage{amsmath,amssymb,amsthm,mathtools,mathrsfs}
\usepackage{booktabs,array,longtable}
\usepackage{enumitem}
\usepackage{xcolor}
\usepackage[hidelinks]{hyperref}
\usepackage[nameinlink,noabbrev]{cleveref}
\hypersetup{
  pdftitle={Residue Transport at the Volume Divisor and Exact Degree-Two Closure in an HSS-Motivated Schoen Sector},
  pdfsubject={Exact Laurent spectral connection and formal-CM rational-function certificate}
}

\newtheorem{theorem}{Theorem}[section]
\newtheorem{proposition}[theorem]{Proposition}
\newtheorem{lemma}[theorem]{Lemma}
\newtheorem{corollary}[theorem]{Corollary}
\theoremstyle{definition}
\newtheorem{definition}[theorem]{Definition}
\newtheorem{certificate}[theorem]{Exact certificate}
\theoremstyle{remark}
\newtheorem{remark}[theorem]{Remark}
\newtheorem{caution}[theorem]{Scope caution}

\newcommand{\dd}{\mathrm{d}}
\newcommand{\ii}{\mathrm{i}}
\newcommand{\e}{\mathrm{e}}
\newcommand{\Q}{\mathbb{Q}}
\newcommand{\Kfield}{\mathbb{K}}
\newcommand{\cB}{\mathcal{B}}
\newcommand{\cR}{\mathcal{R}}
\newcommand{\bK}{\overline{K}}
\newcommand{\bg}{\overline{g}}
\newcommand{\bsigma}{\overline{\sigma}}
\newcommand{\bD}{\overline{D}}
\newcommand{\bSc}{\overline{\mathcal{S}}_{\mathrm{curv}}}
\newcommand{\rank}{\operatorname{rank}}
\newcommand{\ord}{\operatorname{ord}}
\newcommand{\supp}{\operatorname{supp}}
\newcommand{\Span}{\operatorname{span}}

\title{\textbf{Residue Transport at the Volume Divisor\\and Exact Degree-Two Closure\\in an HSS-Motivated Schoen Sector}\\[2mm]
\large Part III: From Constant-Coefficient Obstruction to a Finite Laurent Connection}
% Insert the author's name before public release.
% \author{Author Name}
\author{}
\date{14 July 2026}

\begin{document}
\maketitle

\begin{abstract}
A preceding square-lattice calculation established an exact leading spectral--curvature closure in a three-dimensional Hessian model motivated by the one-sectional Hosono--Saito--Stienstra sector of Schoen's Calabi--Yau threefold.  A second study found a stable degree-two obstruction when the same two response channels were assigned constant correction coefficients, even after canonical holomorphic-anomaly sources, quasi-Jacobi ambiguities, logarithmic spectral terms, contact jets, and second-order operator deformations were added.

Here we resolve that obstruction.  On the symmetric line, the intrinsic coordinate
\[
 U=3X+1
\]
is simultaneously the classical-volume coordinate, the degeneration divisor of the background Hessian metric, and the determinant of the symmetry-adapted frame up to a constant.  The two leading response channels are finite Laurent polynomials in $U$.  We prove that the complete factorized degree-two defect
\[
 R_{\mathrm{quad}}+a\,r_A+b\,r_B+c\,r_C
\]
closes exactly on these same two channels after their degree-two coefficients are promoted from constants to a finite Laurent connection.  The polar coefficients are uniquely determined by residue matching at $U=0$, while the positive-power tail is universal and lies entirely in the $\bsigma$ channel.  The deepest residues satisfy
\[
 \beta_{-3}=\rho_3\alpha_{-3},
 \qquad
 \rho_3=\frac{2(L^2E_4(i)+108)}{3(L^2E_4(i)+36)}.
\]
The universal drift is
\[
 \alpha_{\mathrm{univ}}(U)
 =\frac{\alpha_0C_0^2}{243}U\left(\frac3L-U\right),
 \qquad
 \beta_{\mathrm{univ}}(U)=0.
\]
The final result is certified by exact rational-function arithmetic over a formal square-lattice/level-two CM coefficient field.  After clearing denominators, the intercept numerator has degree at most $13$ and the three geometry-slope numerators have degree at most $10$; exact vanishing at $14$ and $11$ rational points, respectively, proves all four identities.  Thus the earlier degree-two obstruction is not a failure of the rank-two response system, but the signature of residue transport over the unique volume/frame divisor.
\end{abstract}

\section{Introduction}

This paper is the third part of a sequence studying a symmetry-reduced Hessian model motivated by the modular enumerative geometry of Schoen's Calabi--Yau threefold.  The first part isolated an exact square-lattice closure at leading order.  The second part showed that the most direct degree-two extension fails if the two closure coefficients remain constant.  The present paper explains that failure and replaces it with an exact positive theorem.

The underlying one-variable series is
\begin{equation}\label{eq:Bdef}
 B(q)=9\prod_{n\ge1}(1-q^n)^{-4}
     =9q^{1/6}\eta(\tau)^{-4},
 \qquad q=\e^{2\pi\ii\tau}.
\end{equation}
At the square lattice $\tau=i$, the leading curvature response closes on two symmetry-adapted spectral channels with constants $\alpha_0,\beta_0$.  When the first genuine degree-two rational-elliptic-surface data were inserted, an exact one-sided jet test and a high-precision factorized Schoen test both obstructed constant-coefficient closure.  Subsequent calculations tested increasingly broad correction spaces: the canonical holomorphic-anomaly source, split anomaly terms, weight-ten quasi-Jacobi ambiguities, scalar and matrix logarithmic spectra, a universal $\bK_1^2/2$ term, formal contact jets, genuine two-variable contact bi-jets, and the full symmetric second-order radial/theta operator envelope.  These enlarged the observable response space but did not reach the defect.

The decisive change was to allow the degree-two coefficients to vary with the radial modulus.  A finite Laurent module in
\[
 U=3X+1
\]
closed on $33$, then $50$, then $70$ independent radial points, while matched modules based on the frame coefficient $2X+1$ failed.  Fixed geometries and arbitrary factorized geometries all selected the same minimal support.  The resulting affine connection was then reconstructed exactly from the original geometry and certified as a rational-function identity.

The main point is therefore not merely that two functions $\alpha_1(X),\beta_1(X)$ can absorb one scalar defect; arbitrary functions would make such a statement vacuous.  The content is that a fixed nine-dimensional Laurent module, reduced by six exact constraints to three geometry-driven master coordinates and two universal drift coefficients, closes identically for arbitrary factorized geometry coefficients $(a,b,c)$.

\subsection*{Status of the results}

\begin{table}[ht]
\centering
\small
\begin{tabular}{p{0.47\textwidth}p{0.18\textwidth}p{0.25\textwidth}}
\toprule
Result & Status & Method \\
\midrule
Classical volume/frame divisor identities & exact & symbolic Hessian calculation \\
Finite Laurent formulas for $\bsigma$ and $\bD$ & exact & direct reconstruction from the leading geometry \\
Deepest-pole relation $\beta_{-3}=\rho_3\alpha_{-3}$ & exact & Laurent cancellation \\
Rank-seven residue recursion and three-master reduction & exact & symbolic linear algebra \\
Universal positive-$U$ drift & exact & formal-CM Laurent solve \\
Affine connection for arbitrary $(a,b,c)$ & exact & simultaneous four-column solve \\
Global rational identities & exact & denominator clearing and degree-bounded exact evaluation \\
Interpretation as the full Schoen degree-two potential & not claimed & factorized sector only \\
\bottomrule
\end{tabular}
\caption{Logical status of the principal statements.}
\label{tab:status}
\end{table}

\section{The Hessian model and the factorized degree-two family}

\subsection{Background cubic and symmetric frame}

Let
\begin{equation}\label{eq:V}
 V(v_0,v_1,v_2)
 =9v_0v_1v_2+\frac32v_1^2v_2+\frac32v_1v_2^2,
\end{equation}
and set
\begin{equation}\label{eq:Lpq}
 L=2\pi,
 \qquad p=\e^{-Lv_0},
 \qquad q_j=\e^{-Lv_j}\quad(j=1,2).
\end{equation}
The background potential is $K_0=-\log V$.  Restrict to the symmetric line
\begin{equation}\label{eq:symmetricline}
 v=(X,1,1),
\end{equation}
and define
\begin{equation}\label{eq:Udef}
 U=3X+1.
\end{equation}
The symmetry-adapted directions are
\begin{equation}\label{eq:frame}
 e=(-(2X+1),1,1),
 \qquad o=(0,1,-1),
 \qquad r=(X,1,1).
\end{equation}

The leading correction has the form
\begin{equation}\label{eq:K1}
 K_1=p\,\bK_1,
\end{equation}
where
\begin{equation}\label{eq:K1bar}
 \bK_1=
 \frac{B_1B_2(1+Lv_0)+Lv_1(\theta B)_1B_2+Lv_2B_1(\theta B)_2}
 {2L^3V}.
\end{equation}
Here $B_j=B(q_j)$ and $\theta=q\,\dd/\dd q$.  The common factor $p$ is handled by the twisted derivatives
\begin{equation}\label{eq:twisted}
 \mathscr D_0=\partial_{v_0}-L,
 \qquad
 \mathscr D_1=\partial_{v_1},
 \qquad
 \mathscr D_2=\partial_{v_2}.
\end{equation}
Let $M_{ij}=\partial_i\partial_jV$.  The first-order metric response is
\[
 \bg^{(1)}_{ij}=\mathscr D_i\mathscr D_j\bK_1.
\]
Define
\begin{equation}\label{eq:mu}
 \mu_e=\frac{\bg^{(1)}(e,e)}{M(e,e)},
 \qquad
 \mu_o=\frac{\bg^{(1)}(o,o)}{M(o,o)},
 \qquad
 \mu_r=\frac{\bg^{(1)}(r,r)}{M(r,r)},
\end{equation}
and the two response channels
\begin{equation}\label{eq:channels}
 \bsigma=\mu_o-\mu_e,
 \qquad
 \bD=\frac{\mu_e+\mu_o}{2}+2\mu_r.
\end{equation}

At the square lattice, put
\begin{equation}\label{eq:C0Y}
 C_0=B(\e^{-2\pi}),
 \qquad
 Y=L^2E_4(i).
\end{equation}
The leading closure coefficients are
\begin{equation}\label{eq:alpha0beta0}
 \alpha_0=\frac{Y+12}{72},
 \qquad
 \beta_0=\frac{-7Y^2+552Y+13392}{288(Y+36)}.
\end{equation}

\subsection{Degree-two rational-surface data}

Set
\begin{equation}\label{eq:Z2P2}
 Z_2(q)=\frac{E_2(q)}{72}B(q)^2,
 \qquad
 P_2(q)=Z_2(q)-\frac18B(q^2).
\end{equation}
The natural factorized degree-two functions are
\begin{align}
 F_A(q_1,q_2)&=P_2(q_1)P_2(q_2),\label{eq:FA}\\
 F_B(q_1,q_2)&=B(q_1^2)B(q_2^2),\label{eq:FB}\\
 F_C(q_1,q_2)&=P_2(q_1)B(q_2^2)+B(q_1^2)P_2(q_2).\label{eq:FC}
\end{align}
For a symmetric function $F_2(q_1,q_2)$, the degree-two HSS coefficientwise operator used in this model is
\begin{equation}\label{eq:K2operator}
 \bK_2[F_2]
 =\frac{F_2(1+2Lv_0)+Lv_1\theta_1F_2+Lv_2\theta_2F_2}{2L^3V}.
\end{equation}
The full order-$p^2$ curvature response contains the direct contribution of $\bK_2$ and every quadratic $K_1\times K_1$ contribution from the inverse metric, curvature numerator, and bisectional denominator.

Write $R_{\mathrm{quad}}$ for the pure quadratic $K_1\times K_1$ contribution.  For $J=A,B,C$, let $r_J$ be the corrected response of $\bK_2[F_J]$ after subtracting the leading-channel terms $\alpha_0\sigma_2+\beta_0D_2$.  The general factorized degree-two defect is
\begin{equation}\label{eq:defect}
 \mathscr F_2(X;a,b,c)
 =R_{\mathrm{quad}}(X)+a\,r_A(X)+b\,r_B(X)+c\,r_C(X).
\end{equation}
The preceding obstruction paper proved that no constant degree-two pair $\alpha_1,\beta_1$ closes \eqref{eq:defect} throughout the radial line within this family.

\section{The unique volume/frame divisor}

\begin{lemma}[Classical orthogonal frame]\label{lem:frame}
Along the symmetric line, the frame \eqref{eq:frame} is orthogonal for the background Hessian metric $g_0=\nabla^2(-\log V)$ and has constant squared norms
\begin{equation}\label{eq:g0norms}
 g_0(e,e)=6,
 \qquad
 g_0(o,o)=2,
 \qquad
 g_0(r,r)=3.
\end{equation}
Moreover,
\begin{equation}\label{eq:cross}
 g_0(e,o)=g_0(e,r)=g_0(o,r)=0.
\end{equation}
\end{lemma}

\begin{proof}
Differentiate \eqref{eq:V}, restrict to $(X,1,1)$, and contract the resulting Hessian with the three vectors.  The accompanying standalone Sage certificate verifies each rational identity by numerator reduction in the exact fraction field.
\end{proof}

\begin{proposition}[Coincident degeneration loci]\label{prop:divisor}
The coordinate $U=3X+1$ satisfies
\begin{equation}\label{eq:divisoridentities}
 V(X,1,1)=3U,
 \qquad
 \det g_0=\frac9{U^2},
 \qquad
 \det(e,o,r)=-2U,
\end{equation}
and
\begin{equation}\label{eq:er}
 e-r=-U(1,0,0).
\end{equation}
Thus $U=0$ is simultaneously the classical-volume divisor, the background metric degeneration divisor, and the frame-collapse divisor.
\end{proposition}

\begin{proof}
The first identity follows directly from \eqref{eq:V}.  The determinant identities follow from the exact Hessian and frame matrices.  Equation \eqref{eq:er} is immediate from \eqref{eq:frame}.
\end{proof}

\begin{remark}
The factor $2X+1$ occurs in the coordinate expression for $e$, but it is not a degeneration divisor.  This explains why Laurent modules based on $U=3X+1$ succeeded in the numerical discovery stage while matched modules based on $2X+1$ failed.
\end{remark}

\section{Finite Laurent response channels}

\begin{theorem}[Exact finite channel formulas]\label{thm:channels}
At the square lattice, the two leading response channels are finite Laurent polynomials in $U$:
\begin{equation}\label{eq:sigmafinite}
 \bsigma(U)=
 \frac{C_0^2\left(
 8L^3U^3-24L^2U^2-3L(Y+84)U-6(Y+108)
 \right)}{5832L^3U^2},
\end{equation}
and
\begin{equation}\label{eq:Dfinite}
 \bD(U)=
 \frac{C_0^2(Y+36)(LU+3)}{1944L^3U^2}.
\end{equation}
Consequently,
\begin{equation}\label{eq:supportchannels}
 \supp(\bsigma)=\{-2,-1,0,1\},
 \qquad
 \supp(\bD)=\{-2,-1\}.
\end{equation}
\end{theorem}

\begin{proof}
The exact program reconstructs $\bK_1$, differentiates it with the twisted derivatives \eqref{eq:twisted}, contracts with the frame, substitutes $X=(U-1)/3$, and reduces the resulting rational functions.  Expanding \eqref{eq:sigmafinite}--\eqref{eq:Dfinite} gives
\begin{align}
 s_{-2}&=-\frac{C_0^2(Y+108)}{972L^3},
 &d_{-2}&=\frac{C_0^2(Y+36)}{648L^3},\label{eq:leadingcoeffs}\\
 s_{-1}&=-\frac{C_0^2(Y+84)}{1944L^2},
 &d_{-1}&=\frac{C_0^2(Y+36)}{1944L^2},\label{eq:nextcoeffs}\\
 s_0&=-\frac{C_0^2}{243L},
 &d_0&=0,\label{eq:constantcoeffs}\\
 s_1&=\frac{C_0^2}{729},
 &d_j&=0\quad(j\ge0),\label{eq:positivecoeffs}
\end{align}
with $s_j=0$ for $j\ge2$.
\end{proof}

\begin{corollary}[Support forcing]\label{cor:supportforcing}
Within the minimal connection support used below, the $\bD$ channel cannot contribute to the regular or positive-power tail.  In particular, every universal positive-$U$ term is forced into the $\bsigma$ channel.
\end{corollary}

\section{The Laurent connection and exact closure}

\begin{definition}[Degree-two Laurent connection]\label{def:connection}
Let
\begin{align}
 \alpha_1(U;a,b,c)
 &={A_{-3}}U^{-3}+{A_{-2}}U^{-2}+{A_{-1}}U^{-1}+A_0+A_1U+A_2U^2,
 \label{eq:alphaansatz}\\
 \beta_1(U;a,b,c)
 &={B_{-3}}U^{-3}+{B_{-2}}U^{-2}+B_0.
 \label{eq:betaansatz}
\end{align}
The coefficients $A_{-3},A_{-2},A_{-1},A_0,B_{-3},B_{-2},B_0$ are allowed to depend affinely on $(a,b,c)$.  The two positive-power coefficients are fixed by
\begin{equation}\label{eq:universalcoeffs}
 A_1=\frac{\alpha_0C_0^2}{81L},
 \qquad
 A_2=-\frac{\alpha_0C_0^2}{243}.
\end{equation}
\end{definition}

\begin{theorem}[Exact degree-two residue-transport closure]\label{thm:main}
Let $\Kfield$ be the rational function field
\begin{equation}\label{eq:field}
 \Kfield=
 \Q\bigl(
 L,E_4(i),C_0,B(\e^{-4\pi}),E_2(2i),E_4(2i),E_6(2i)
 \bigr),
\end{equation}
with the denominators occurring below inverted.  Generate every required theta jet from Ramanujan's differential equations and
\[
 \theta B=\frac{1-E_2}{6}B.
\]
For the exact degree-two defect \eqref{eq:defect}, there exists a unique connection of the form \eqref{eq:alphaansatz}--\eqref{eq:betaansatz}, with the drift \eqref{eq:universalcoeffs}, such that
\begin{equation}\label{eq:mainidentity}
 \boxed{
 \mathscr F_2(X;a,b,c)
 =\alpha_1(U;a,b,c)\,\bsigma(U)
 +\beta_1(U;a,b,c)\,\bD(U)
 }
\end{equation}
identically as a rational function of $X$, for arbitrary $(a,b,c)$.

The seven non-universal coefficients depend affinely on $(a,b,c)$.  Their geometry-slope map has rank three.  Equivalently, they are determined by three master coordinates, which may be chosen as
\begin{equation}\label{eq:masters}
 A_0,
 \qquad B_0,
 \qquad A_{-3}.
\end{equation}
\end{theorem}

\begin{proof}
The proof has four exact stages.

First, Ramanujan's equations generate the square-lattice and level-two theta jets through the required order.  The level-two values in \eqref{eq:field} remain formal; hence the calculation is stronger than a single numerical specialization.

Second, a truncated multivariate Taylor-jet algebra reconstructs $K_0$, $K_1$, the full quadratic $K_1\times K_1$ curvature response, and the direct degree-two responses of $F_A,F_B,F_C$.  This produces exact rational functions $R_{\mathrm{quad}},r_A,r_B,r_C,\bsigma,\bD$.

Third, Laurent expansion at $U=0$ gives a rank-seven linear system for the seven non-universal coefficients in each of the four affine columns: the intercept, the $a$ slope, the $b$ slope, and the $c$ slope.  The same coefficient matrix is invertible by \cref{prop:polar}.  Solving the four right-hand sides simultaneously yields an exact affine connection map.

Fourth, each of the four identities is cleared of denominators.  After multiplying the connection by $U^3$, the intercept cleared numerator has degree at most $13$, while the three slope numerators have degree at most $10$.  Exact evaluation at $14$ and $11$ distinct rational points, respectively, proves that the four polynomials vanish identically.  The certificate results are recorded in \cref{tab:certificate}.  Affine linearity then proves \eqref{eq:mainidentity} for arbitrary $(a,b,c)$.
\end{proof}

\begin{corollary}[Universal drift]\label{cor:drift}
The geometry-independent part of the connection is
\begin{equation}\label{eq:drift}
 \boxed{
 \alpha_{\mathrm{univ}}(U)
 =\frac{\alpha_0C_0^2}{243}U\left(\frac3L-U\right),
 \qquad
 \beta_{\mathrm{univ}}(U)=0.
 }
\end{equation}
\end{corollary}

\section{Residue recursion at the frame divisor}

Write
\[
 \bsigma=\sum_j s_jU^j,
 \qquad
 \bD=\sum_j d_jU^j.
\]
The nominal deepest term in \eqref{eq:mainidentity} is $U^{-5}$.

\begin{proposition}[Deepest-pole alignment]\label{prop:rho}
The absence of a $U^{-5}$ defect term forces
\begin{equation}\label{eq:rhorelation}
 B_{-3}=\rho_3A_{-3},
\end{equation}
where
\begin{equation}\label{eq:rho}
 \boxed{
 \rho_3=-\frac{s_{-2}}{d_{-2}}
 =\frac{2(Y+108)}{3(Y+36)}
 =\frac{2(E_4+3E_2^2)}{3(E_4+E_2^2)}\bigg|_{\tau=i}.
 }
\end{equation}
\end{proposition}

\begin{proof}
The coefficient of $U^{-5}$ is
\[
 A_{-3}s_{-2}+B_{-3}d_{-2}.
\]
Set it to zero and use \eqref{eq:leadingcoeffs}.  The last expression follows from $LE_2(i)=6$.
\end{proof}

The next three equations determine $A_{-2},B_{-2},A_{-1}$.  Their coefficient matrix is
\begin{equation}\label{eq:polarmatrix}
 \mathsf M_{\mathrm{pol}}=
 \begin{pmatrix}
 s_{-2}&d_{-2}&0\\
 s_{-1}&d_{-1}&s_{-2}\\
 s_0&0&s_{-1}
 \end{pmatrix}.
\end{equation}

\begin{proposition}[Polar nondegeneracy]\label{prop:polar}
The determinant of \eqref{eq:polarmatrix} is
\begin{equation}\label{eq:polardet}
 \det\mathsf M_{\mathrm{pol}}
 =-\frac{C_0^6(Y+36)(Y^2+72Y-2160)}{7346640384L^7}.
\end{equation}
It is nonzero at the physical square-lattice point.  Hence the polar recursion is unique.
\end{proposition}

\begin{remark}[Actual valuations]
The exact reconstructed geometry gives
\begin{equation}\label{eq:valuations}
 \ord_{U=0}R_{\mathrm{quad}}=-3,
 \qquad
 \ord_{U=0}r_A
 =\ord_{U=0}r_B
 =\ord_{U=0}r_C=-2.
\end{equation}
Thus both the $U^{-5}$ and $U^{-4}$ response terms cancel before the first intercept residue appears.  The geometry slopes begin one order later still.
\end{remark}

\section{Finite support and the universal polynomial tail}

By \cref{thm:channels}, the intercept connection product has possible support from $U^{-5}$ to $U^3$.  The exact cancellations and valuations sharpen this.

\begin{corollary}[Finite Laurent support]\label{cor:support}
The exact intercept defect has support contained in $[-3,3]$.  Each geometry slope has support contained in $[-2,1]$.  Therefore the complete affine defect \eqref{eq:defect} is a finite Laurent polynomial in $U$.
\end{corollary}

The top two intercept coefficients contain no geometry-dependent master parameters.

\begin{proposition}[Universal top coefficients]\label{prop:top}
Let
\[
 \mathscr F_2(U;a,b,c)=\sum_kF_k(a,b,c)U^k.
\]
Then
\begin{equation}\label{eq:F3}
 \boxed{F_3=-\frac{\alpha_0C_0^4}{3^{11}}},
\end{equation}
and
\begin{equation}\label{eq:F2}
 \boxed{F_2=\frac{2\alpha_0C_0^4}{3^{10}L}}.
\end{equation}
Both are independent of $(a,b,c)$.
\end{proposition}

\begin{proof}
Only $A_2U^2\cdot s_1U$ contributes to $F_3$, while
\[
 F_2=A_1s_1+A_2s_0.
\]
Insert \eqref{eq:universalcoeffs}, \eqref{eq:constantcoeffs}, and \eqref{eq:positivecoeffs}.
\end{proof}

\section{The exact formal-CM certificate}

\subsection{Coefficient field and generated jets}

The exact program works over the formal field \eqref{eq:field}.  Let
\[
 E_k^{(2)}=E_k(2i),
 \qquad
 B^{(2)}=B(\e^{-4\pi}).
\]
The required higher derivatives are not introduced as independent variables.  They are generated recursively from Ramanujan's equations
\begin{equation}\label{eq:Ramanujan}
 \theta E_2=\frac{E_2^2-E_4}{12},
 \qquad
 \theta E_4=\frac{E_2E_4-E_6}{3},
 \qquad
 \theta E_6=\frac{E_2E_6-E_4^2}{2},
\end{equation}
and
\begin{equation}\label{eq:Bderiv}
 \theta B=\frac{1-E_2}{6}B.
\end{equation}
For a doubled argument,
\begin{equation}\label{eq:doublejet}
 \theta_q^m f(q^2)=2^m(\theta^mf)(q^2).
\end{equation}
This produces exact jets of $B$, $Z_2$, $B(q^2)$, and $P_2$ through the required order.

\subsection{Degree-bounded polynomial proof}

Writing
\[
 F=\frac{n_F}{d_F},
 \qquad
 \bsigma=\frac{n_\sigma}{d_\sigma},
 \qquad
 \bD=\frac{n_D}{d_D},
\]
and defining
\[
 \widetilde A=U^3\alpha_1,
 \qquad
 \widetilde B=U^3\beta_1,
\]
the cleared numerator is
\begin{equation}\label{eq:clearedN}
 \begin{aligned}
 N(X)={}&n_Fd_\sigma d_DU^3\\
 &-d_F\left(
 \widetilde A\,n_\sigma d_D
 +\widetilde B\,n_Dd_\sigma
 \right).
 \end{aligned}
\end{equation}
This is a polynomial in $X$ over $\Kfield$.  Rather than globally expanding its enormous formal-CM coefficients, the verifier computes a rigorous degree bound and evaluates \eqref{eq:clearedN} exactly at sufficiently many rational points.

\begin{certificate}[Final exact certificate]\label{cert:final}
The four affine columns produced the following results.
\begin{table}[ht]
\centering
\begin{tabular}{lccc}
\toprule
Column & Degree bound & Exact evaluations & Certified zero \\
\midrule
Intercept $R_{\mathrm{quad}}$ & $13$ & $14$ & true \\
Slope $r_A$ & $10$ & $11$ & true \\
Slope $r_B$ & $10$ & $11$ & true \\
Slope $r_C$ & $10$ & $11$ & true \\
\bottomrule
\end{tabular}
\caption{Exact cleared-polynomial certificates.  Every evaluation was performed in the formal coefficient field, with no floating-point tolerance.}
\label{tab:certificate}
\end{table}
A polynomial of degree at most $D$ that vanishes at $D+1$ distinct points is zero.  Hence all four rational identities hold identically.
\end{certificate}

\section{How the obstruction became a theorem}

The degree-two obstruction in the companion paper was real but conditional on constant coefficients.  It persisted under many natural potential and operator enlargements because those experiments searched for a new fixed response direction.  The exact result here shows that the missing structure is instead transport within the original rank-two response bundle.

This interpretation explains several empirical facts from the discovery stage.

\begin{enumerate}[label=\arabic*.]
 \item Pure one-variable contact jets were redundant modulo the existing response span.
 \item Genuine mixed contact jets and second-order operator terms added observable directions, but the defect remained independent of them.
 \item A Laurent module in $U=3X+1$ closed, while a matched module in $2X+1$ did not.
 \item The same minimal support worked for the primitive product, primitive-plus-double-cover, $Z_2\times Z_2$, and arbitrary factorized geometry.
 \item The positive-power coefficients were geometry-independent from the outset.
\end{enumerate}

The final theorem does not say that no possible constant-channel enlargement could close.  It says that, for the stated factorized degree-two model, a finite Laurent connection on the original two channels gives an exact and geometrically natural closure, with poles controlled by the unique frame divisor.

\section{Scope and limitations}

\begin{caution}
The exact theorem concerns the factorized family \eqref{eq:FA}--\eqref{eq:FC} in the stated three-dimensional HSS-motivated Hessian model.  It is not asserted that this family is the complete base-degree-two relative Schoen potential.
\end{caution}

More specifically:

\begin{enumerate}[label=\arabic*.]
 \item The identity holds for arbitrary coefficients $(a,b,c)$ multiplying the three factorized components, but no claim is made that a particular triple is the unique physical Schoen coefficient vector.
 \item The formal level-two CM field keeps $B(\e^{-4\pi})$, $E_2(2i)$, $E_4(2i)$, and $E_6(2i)$ symbolic while imposing their Ramanujan differential jets.  This proves an identity stronger than numerical specialization, but it does not by itself identify the smallest algebraic CM field.
 \item The theorem concerns the symmetric radial slice and the even--odd primitive curvature response, not the full special-K\"ahler tensor on the complete moduli space.
 \item The term ``connection'' refers to the finite Laurent transport law for the two response coefficients.  No claim is made here that it is the restriction of a globally defined flat connection on the full moduli space.
\end{enumerate}

\section{Reproducibility and supplementary files}

The accompanying bundle contains the LaTeX source, the compiled PDF, a certificate summary, and the SageMath programs used in the discovery and exact proof.  The principal files are listed below.

\begin{description}[style=nextline,leftmargin=0pt,labelsep=0pt,itemsep=4pt]
\small
\item[\path{22_evolving_spectral_connection_test.sage}]
First finite-Laurent closure on $33$ radial points; $U$ positive control and $2X+1$ negative control.
\item[\path{23_fixed_geometry_U_connection_test.sage}]
Fixed-geometry tests on $50$ points and minimal support search.
\item[\path{24_affine_geometry_connection_reconstruction.sage}]
Affine geometry-to-connection map, $70$-point validation, and universal coefficients.
\item[\path{25_connection_constraint_and_universal_drift_restart_safe_v2.sage}]
Rank-three slope map, six affine constraints, and universal drift recognition.
\item[\path{26_volume_divisor_laurent_matching_standalone_v3.sage}]
Standalone exact frame divisor, finite channel data, $\rho_3$, and polar recursion.
\item[\path{27_exact_degree2_defect_certificate.sage}]
Monolithic exact formal-CM geometry and rational-identity calculation.
\item[\path{27_exact_degree2_cocalc_checkpointed_v2.sage}]
Checkpointed CoCalc version through the affine Laurent solve.
\item[\path{27c_fast_verify_from_checkpoint_v2.sage}]
Degree-bounded exact verifier loaded from the Part-IV checkpoint.
\item[\path{27d_verify_*_restartable.sage}]
One-column restartable exact verifiers for the intercept and three slopes.
\item[\path{27e_aggregate_column_certificates.sage}]
Aggregates the four exact column certificates.
\end{description}

The final certificate summary is
\begin{verbatim}
EXACT DEGREE-TWO FAST CERTIFICATE
intercept Rquad: True (degree <= 13)
slope rA: True (degree <= 10)
slope rB: True (degree <= 10)
slope rC: True (degree <= 10)
\end{verbatim}

\section{Conclusion}

The leading square-lattice closure survives degree two, but not as a constant-coefficient identity.  The correct extension is a finite Laurent connection over the unique volume/frame coordinate $U=3X+1$.  Its polar sector is fixed by residue cancellation at the frame divisor, its geometry dependence has rank three, and its polynomial tail is universal and supported entirely in the $\bsigma$ channel.  The full factorized defect closes exactly for arbitrary $(a,b,c)$ over the formal CM coefficient field.

The conceptual progression across the three papers is therefore:
\begin{equation*}
\begin{gathered}
 \text{CM fixed-point closure}
 \quad\longrightarrow\quad
 \text{constant-coefficient degree-two obstruction}\\
 \quad\longrightarrow\quad
 \text{exact residue transport at the volume divisor}.
\end{gathered}
\end{equation*}

\appendix

\section{Compact square-lattice constants}

The standard square-lattice evaluation gives
\begin{equation}\label{eq:E4gamma}
 E_4(i)=\frac{3\Gamma(1/4)^8}{(2\pi)^6}.
\end{equation}
Moreover,
\begin{equation}\label{eq:C0gamma}
 C_0=B(\e^{-2\pi})
 =\frac{144\pi^3\e^{-\pi/3}}{\Gamma(1/4)^4}.
\end{equation}
Hence the universal drift coefficients may be written
\begin{equation}\label{eq:A1gamma}
 A_1=\frac{128\alpha_0\pi^5\e^{-2\pi/3}}{\Gamma(1/4)^8},
\end{equation}
and
\begin{equation}\label{eq:A2gamma}
 A_2=-\frac{256\alpha_0\pi^6\e^{-2\pi/3}}{3\Gamma(1/4)^8}.
\end{equation}
Their ratio is $A_2/A_1=-L/3$.

\section{Symmetry and the even--radial block}

Under the involution $v_1\leftrightarrow v_2$, the vectors $e,r$ are even and $o$ is odd.  Every exchange-symmetric correction therefore satisfies
\[
 g_j(e,o)=g_j(r,o)=0.
\]
The odd direction is an exact one-dimensional invariant block, while the even and radial directions may mix.  This symmetry explains the matrix-block structure encountered in the obstruction calculations preceding the Laurent-connection discovery.

\section{Certificate manifest}

The bundle also includes a machine-readable manifest and a README containing recommended execution paths for SageCell and CoCalc.  Checkpoint files created at runtime are not included because they are machine- and session-dependent; the supplied scripts reproduce them.

\begin{thebibliography}{99}

\bibitem{HSS}
S.~Hosono, M.-H.~Saito, and J.~Stienstra,
\newblock On the mirror symmetry conjecture for Schoen's Calabi--Yau 3-folds,
\newblock in \emph{Integrable Systems and Algebraic Geometry}, World Scientific, 1998, pp.~194--235;
\newblock arXiv:alg-geom/9709027.

\bibitem{HST}
S.~Hosono, M.-H.~Saito, and A.~Takahashi,
\newblock Holomorphic anomaly equation and BPS state counting of rational elliptic surface,
\newblock \emph{Advances in Theoretical and Mathematical Physics} \textbf{3} (1999), 177--208;
\newblock arXiv:hep-th/9901151.

\bibitem{OP}
G.~Oberdieck and A.~Pixton,
\newblock Gromov--Witten theory of elliptic fibrations: Jacobi forms and holomorphic anomaly equations,
\newblock \emph{Geometry \& Topology} \textbf{23} (2019), 1415--1489;
\newblock arXiv:1709.01481.

\bibitem{Schoen}
C.~Schoen,
\newblock On fiber products of rational elliptic surfaces with section,
\newblock \emph{Mathematische Zeitschrift} \textbf{197} (1988), 177--199.

\bibitem{Ramanujan}
S.~Ramanujan,
\newblock On certain arithmetical functions,
\newblock \emph{Transactions of the Cambridge Philosophical Society} \textbf{22} (1916), 159--184.

\bibitem{Zagier}
D.~Zagier,
\newblock Elliptic modular forms and their applications,
\newblock in \emph{The 1-2-3 of Modular Forms}, Springer, 2008, pp.~1--103.

\bibitem{Sage}
The Sage Developers,
\newblock \emph{SageMath, the Sage Mathematics Software System},
\newblock open-source mathematical software.

\end{thebibliography}

\end{document}
